% GENERATED FILE. Edit canonical lessons, then run npm run generate:books.
% canonical-source-hash: fnv1a64:fe8c17f7e6f35d9e
% canonical-lessons: LA-C345-rex, LA-C345-regno, LA-C345-tyrannus, LA-C345-plebs, LA-C345-dictator

\chapter{Rēgēs et dictātōrēs}
\label{ch:la-reges-et-dictatores}
\hlchaptermodality{345}

\begin{chapteropening}
\textbf{By the end of this chapter:} \emph{I can talk about early rulers and the people: a king, to reign, a tyrant, the common people and a dictator.}

Complete the last lesson of chapter 345: I can talk about early rulers and the people: a king, to reign, a tyrant, the common people and a dictator.
\end{chapteropening}

\section[rēx, rēgis]{rēx, rēgis — a king}

\label{lesson:LA-C345-rex}



\begin{quote}
Before the new one: say the Latin for an origin, then the Latin for great age.
\end{quote}

\subsection*{You'll want to know: rēx, rēgis}
\textbf{rēx, rēgis} — \textquotedblleft{}a king\textquotedblright{}.

\begin{quote}
Ōlim rēgēs hanc urbem tenēbant.
\end{quote}

\begin{quote}
\emph{Once kings held this city.}
\end{quote}

Masculine, 3rd declension. After the last king was driven out, the Romans hated the very name \emph{rēx}.

\begin{cousinweb}
\textbf{rēx} is inherited from Proto-Indo-European \emph{*h\textsubscript{3}r\'{\={e}}ǵs}, \textquotedblleft{}ruler, king\textquotedblright{} (EtymDB), a cousin of Sanskrit \emph{rāj-}. English \emph{regal} and \emph{royal} come from its family.
\end{cousinweb}

\begin{grammarlens}[title={What you've built}]
One more word for telling, explaining and describing at B1.
\end{grammarlens}

\subsection*{Your turn}
\begin{itemize}
\raggedright
  \item \emph{Say it:} \emph{rēx, rēgis}
  \item \emph{Say it:} \emph{rēx, rēgis}, once more
  \item \emph{Say it:} a sentence of your own with \emph{rēx, rēgis}
  \item \emph{Recall:} say the Latin for an origin, then the Latin for great age, then say \emph{rēx, rēgis} again
\end{itemize}

\begin{tcolorbox}[breakable,colback=teal!4,colframe=teal!35!black,title={Before you move on}]
What does rēx, rēgis mean? (\textquotedblleft{}A king\textquotedblright{}.) Say it once more.
\end{tcolorbox}
\section[rēgnō, rēgnāre]{rēgnō, rēgnāre — to reign, to be king}

\label{lesson:LA-C345-regno}



\begin{quote}
Before the new one: say the Latin for great age, then the Latin for a king.
\end{quote}

\subsection*{You'll want to know: rēgnō, rēgnāre}
\textbf{rēgnō, rēgnāre} — \textquotedblleft{}to reign, to be king\textquotedblright{}.

\begin{quote}
Senex rēx multōs annōs rēgnāvit.
\end{quote}

\begin{quote}
\emph{The old king reigned for many years.}
\end{quote}

\emph{rēgnō, rēgnāre, rēgnāvī, rēgnātum}, from \emph{rēgnum} (taught). English \emph{reign}.

\begin{cousinweb}
\textbf{rēgnō} is built on \emph{rēgnum}, \textquotedblleft{}kingship, kingdom\textquotedblright{} (EtymDB). English \emph{reign} comes from its family through French.
\end{cousinweb}

\begin{grammarlens}[title={What you've built}]
One more word for telling, explaining and describing at B1.
\end{grammarlens}

\subsection*{Your turn}
\begin{itemize}
\raggedright
  \item \emph{Say it:} \emph{rēgnō, rēgnāre}
  \item \emph{Say it:} \emph{rēgnō, rēgnāre}, once more
  \item \emph{Say it:} a sentence of your own with \emph{rēgnō, rēgnāre}
  \item \emph{Recall:} say the Latin for great age, then the Latin for a king, then say \emph{rēgnō, rēgnāre} again
\end{itemize}

\begin{tcolorbox}[breakable,colback=teal!4,colframe=teal!35!black,title={Before you move on}]
What does rēgnō, rēgnāre mean? (\textquotedblleft{}To reign, to be king\textquotedblright{}.) Say it once more.
\end{tcolorbox}
\section[tyrannus, tyrannī]{tyrannus, tyrannī — a tyrant, a despot}

\label{lesson:LA-C345-tyrannus}



\begin{quote}
Before the new one: say the Latin for a king, then the Latin for to reign.
\end{quote}

\subsection*{You'll want to know: tyrannus, tyrannī}
\textbf{tyrannus, tyrannī} — \textquotedblleft{}a tyrant, a despot\textquotedblright{}.

\begin{quote}
Cīvēs tyrannum ex urbe pepulērunt.
\end{quote}

\begin{quote}
\emph{The citizens drove the tyrant out of the city.}
\end{quote}

Masculine. For a Roman, a ruler who seized power and ruled without law.

\begin{cousinweb}
\textbf{tyrannus} was borrowed from Greek \emph{tyrannos}, \textquotedblleft{}absolute ruler\textquotedblright{} (EtymDB). English \emph{tyrant} comes from it through French.
\end{cousinweb}

\begin{grammarlens}[title={What you've built}]
One more word for telling, explaining and describing at B1.
\end{grammarlens}

\subsection*{Your turn}
\begin{itemize}
\raggedright
  \item \emph{Say it:} \emph{tyrannus, tyrannī}
  \item \emph{Say it:} \emph{tyrannus, tyrannī}, once more
  \item \emph{Say it:} a sentence of your own with \emph{tyrannus, tyrannī}
  \item \emph{Recall:} say the Latin for a king, then the Latin for to reign, then say \emph{tyrannus, tyrannī} again
\end{itemize}

\begin{tcolorbox}[breakable,colback=teal!4,colframe=teal!35!black,title={Before you move on}]
What does tyrannus, tyrannī mean? (\textquotedblleft{}A tyrant, a despot\textquotedblright{}.) Say it once more.
\end{tcolorbox}
\section[plēbs, plēbis]{plēbs, plēbis — the common people}

\label{lesson:LA-C345-plebs}



\begin{quote}
Before the new one: say the Latin for to reign, then the Latin for a tyrant.
\end{quote}

\subsection*{You'll want to know: plēbs, plēbis}
\textbf{plēbs, plēbis} — \textquotedblleft{}the common people\textquotedblright{}.

\begin{quote}
Plēbs in forō clāmābat.
\end{quote}

\begin{quote}
\emph{The common people were shouting in the forum.}
\end{quote}

Feminine, 3rd declension, singular. The citizens outside the old noble families. English \emph{plebeian}.

\begin{cousinweb}
\textbf{plēbs} is linked by Lewis \& Short with a root meaning \textquotedblleft{}to fill\textquotedblright{}, as in \emph{plēnus}, \textquotedblleft{}full\textquotedblright{}: the multitude. English \emph{plebeian} comes from its family.
\end{cousinweb}

\begin{grammarlens}[title={What you've built}]
One more word for telling, explaining and describing at B1.
\end{grammarlens}

\subsection*{Your turn}
\begin{itemize}
\raggedright
  \item \emph{Say it:} \emph{plēbs, plēbis}
  \item \emph{Say it:} \emph{plēbs, plēbis}, once more
  \item \emph{Say it:} a sentence of your own with \emph{plēbs, plēbis}
  \item \emph{Recall:} say the Latin for to reign, then the Latin for a tyrant, then say \emph{plēbs, plēbis} again
\end{itemize}

\begin{tcolorbox}[breakable,colback=teal!4,colframe=teal!35!black,title={Before you move on}]
What does plēbs, plēbis mean? (\textquotedblleft{}The common people\textquotedblright{}.) Say it once more.
\end{tcolorbox}
\section[dictātor, dictātōris]{dictātor, dictātōris — a dictator}

\label{lesson:LA-C345-dictator}



\begin{quote}
Before the new one: say the Latin for a tyrant, then the Latin for the common people.
\end{quote}

\subsection*{You'll want to know: dictātor, dictātōris}
\textbf{dictātor, dictātōris} — \textquotedblleft{}a dictator\textquotedblright{}.

\begin{quote}
In magnō perīculō senātus dictātōrem creāvit.
\end{quote}

\begin{quote}
\emph{In great danger the senate appointed a dictator.}
\end{quote}

Masculine, 3rd declension. The Roman dictator held full power in an emergency, for six months only.

\begin{cousinweb}
\textbf{dictātor} is built on \emph{dictō}, \textquotedblleft{}to say often, to order\textquotedblright{} (Lewis \& Short), from \emph{dīcō}, \textquotedblleft{}to say\textquotedblright{}. English \emph{dictator} is the same word.
\end{cousinweb}

\begin{grammarlens}[title={What you've built}]
One more word for telling, explaining and describing at B1.
\end{grammarlens}

\subsection*{Your turn}
\begin{itemize}
\raggedright
  \item \emph{Say it:} \emph{dictātor, dictātōris}
  \item \emph{Say it:} \emph{dictātor, dictātōris}, once more
  \item \emph{Say it:} a sentence of your own with \emph{dictātor, dictātōris}
  \item \emph{Recall:} say the Latin for a tyrant, then the Latin for the common people, then say \emph{dictātor, dictātōris} again
\end{itemize}

\begin{tcolorbox}[breakable,colback=teal!4,colframe=teal!35!black,title={Before you move on}]
What does dictātor, dictātōris mean? (\textquotedblleft{}A dictator\textquotedblright{}.) Say it once more.
\end{tcolorbox}
